% Unit 5 review sheet -- 2 pages.
\documentclass{apchem-worksheet}

\apcunit{5}
\apcunittitle{Kinetics}
\apcdoctitle{Review Sheet}
\apcsubtitle{Everything on the exam traces to this page \textbullet\ exam Fri Nov 20}

\begin{document}
\apctitleblock

\section*{\sffamily The five tools}

\renewcommand{\arraystretch}{1.45}
\noindent\begin{tabular}{@{}p{3.9cm}p{9.8cm}@{}}
\toprule
\textbf{Rate definition} & rate $= -\tfrac{1}{a}\Delta[\ce{A}]/\Delta t =
  +\tfrac{1}{c}\Delta[\ce{C}]/\Delta t$. \emph{Always divide by the
  coefficient} --- otherwise ``the rate'' depends on which species you
  watched\\
\textbf{Rate law} & rate $= k[\ce{A}]^m[\ce{B}]^n$, orders from
  \emph{experiment}. Method of initial rates: change one concentration,
  see what the rate does. Doubling $\to$ $\times1$, $\times2$, $\times4$
  means order 0, 1, 2\\
\textbf{Integrated laws} & zero: $[\ce{A}]$ vs $t$ linear; first:
  $\ln[\ce{A}]$ vs $t$; second: $1/[\ce{A}]$ vs $t$. Slope is $-k$, $-k$,
  $+k$\\
\textbf{Mechanism tests} & steps must sum to the overall equation AND
  predict the observed rate law. Elementary steps only: molecularity gives
  order\\
\textbf{Collision model} & rate depends on collision frequency, energy
  ($\ge E_a$) and orientation. $k = Ae^{-E_a/RT}$\\
\bottomrule
\end{tabular}

\section*{\sffamily Half-lives, and what they tell you}

\begin{itemize}[leftmargin=*,itemsep=0.18em]
  \item \textbf{Zero order:} $t_{1/2} = [\ce{A}]_0/2k$ --- gets
        \emph{shorter} as the reaction proceeds.
  \item \textbf{First order:} $t_{1/2} = 0.693/k$ --- \emph{constant}, the
        only one independent of concentration. All radioactive decay.
  \item \textbf{Second order:} $t_{1/2} = 1/(k[\ce{A}]_0)$ --- gets
        \emph{longer} as the reaction proceeds.
  \item A half-life that does not change is therefore direct evidence of
        first-order kinetics.
\end{itemize}

\section*{\sffamily Mechanisms in five lines}

\begin{itemize}[leftmargin=*,itemsep=0.18em]
  \item \textbf{Slow step first} $\Rightarrow$ the rate law is that step's,
        written straight from its molecularity.
  \item \textbf{Fast step first} $\Rightarrow$ write the slow step's rate
        law, then replace the intermediate using $K$ of step 1.
  \item \textbf{Intermediate:} produced then consumed. Never in a final
        rate law.
  \item \textbf{Catalyst:} consumed then regenerated. Present at the end.
  \item Matching the rate law makes a mechanism \emph{consistent}, never
        \emph{proven}.
\end{itemize}

\section*{\sffamily The traps, one line each}

\begin{itemize}[leftmargin=*,itemsep=0.22em]
  \item Reading orders off the balanced equation.
        ($\ce{2N2O5 -> 4NO2 + O2}$ is \emph{first} order.)
  \item Forgetting the coefficient when converting between species rates.
  \item Saying ``the graph is linear'' without naming the plotted quantity.
  \item Using $t_{1/2} = 0.693/k$ for a reaction that is not first order.
  \item Leaving an intermediate in a final rate law.
  \item Saying a catalyst shifts the equilibrium or changes $\Delta H$.
  \item Reporting $k$ without units --- the units are a free check on your
        orders.
\end{itemize}

\section*{\sffamily Self-check --- 10 questions, exam conditions}

\begin{enumerate}[leftmargin=*,itemsep=0.4em]
  \item For $\ce{2A -> 3B}$, B forms at \SI{6.0e-4}{\Molar\per\second}. How
        fast does A disappear?
        \answerblank[3.4cm]{\SI{4.0e-4}{\Molar\per\second}}
  \item Doubling $[\ce{A}]$ multiplies the rate by 4. Order in A?
        \answerblank[1.8cm]{2}
  \item Units of $k$ for a second-order reaction:
        \answerblank[3.2cm]{\si{\per\Molar\per\second}}
  \item Which plot is linear for first order?
        \answerblank[3.4cm]{$\ln[\ce{A}]$ vs $t$}
  \item $k = \SI{0.0693}{\per\second}$; find $t_{1/2}$:
        \answerblank[2.6cm]{\SI{10.0}{\second}}
  \item Rate law for the elementary step $\ce{2NO -> N2O2}$:
        \answerblank[3.4cm]{rate $= k[\ce{NO}]^2$}
  \item $E_a(\text{fwd}) = 80$, $E_a(\text{rev}) = 120$ kJ. $\Delta H$?
        \answerblank[2.6cm]{$-40$ kJ}
  \item A catalyst changes $K$ how? \answerblank[2.6cm]{not at all}
  \item In $\ce{X + C -> XC}$ then $\ce{XC + Y -> XY + C}$, name the
        catalyst: \answerblank[2cm]{C}
  \item Constant half-life implies which order?
        \answerblank[2.4cm]{first}
\end{enumerate}

\begin{apcnote}[How to study for this exam]
Redo WS 4's FRQ cold --- the exam's long free-response is its twin in shape
and length. Then drill the chain \emph{table $\to$ orders $\to$ rate law
$\to$ $k$ with units} until it runs without stopping, because it opens
almost every kinetics FRQ.

For every ``explain'' item, say what the \emph{particles} are doing before
naming the effect: more collisions, more energetic collisions, a lower
barrier, a different pathway. And whenever a graph is mentioned, name the
axis.
\end{apcnote}

\end{document}
